\section{The charge near separated centers}
\label{fast:section}

Near a commuting configuration with distinct joint centers, the off-block
coordinates form a fast oscillator. We construct its ground-line
projection and prove concentration on this line, together with bounds
for the charge in the block variables. We then
compute the charge compressed to the line at scalar blocks. These local
estimates do not require invariance under the rotation torus $G$.
Gauge equivariance removes the longitudinal fast momentum and the
central gauge angular momentum in the oscillator square.

The fixed-mass regime supplies cluster profiles, while the growing-mass
regime will exclude escape to infinity after rescaling large annuli.
We use two parameter regimes. In the first, $m\geq0$ is fixed, $h\to\infty$,
and a bounded family has locally bounded graph norm for one fixed charge
$D_v(h,m)$. In the second,
\begin{equation}
 m=m_h\longrightarrow\infty,\qquad m_h\leq C_0\sqrt h,
 \qquad \|D_{v_\alpha}(h,m_h)\psi_h\|=o(m_h)
 \quad(\alpha=1,\ldots,16),
 \label{fast:growing-mass}
\end{equation}
where $(v_\alpha)$ is an orthonormal spin basis and
$\sup_h\|\psi_h\|<\infty$. Every assertion with local hypotheses means the
following: multiply by a fixed smooth gauge-invariant cutoff supported in
the specified compact neighborhood, and assume the stated graph bounds
there. Such multiplication preserves the bounds, because the commutator
with a charge is Clifford multiplication by the cutoff gradient and is
independent of $h$ and $m$. All constants below may depend on that
neighborhood and cutoff, $N$, and, in the second regime, $C_0$.

\subsection{Local estimates and spectral slices}

\begin{lemma}[Raw local estimates]\label{lem:raw-estimates}
In either regime, for states supported in a fixed compact set,
\begin{equation}
 \|\nabla\psi_h\|\leq C\sqrt h,
 \qquad
 \sum_{i<j}\|[x_i,x_j]\psi_h\|\leq C h^{-1/2}.
 \label{eq:raw-estimates}
\end{equation}
Consequently every weak limit of the squared norm measures on compact
sets is supported on commuting configurations.
\end{lemma}

\begin{proof}
For fixed $m$, its multiplication term is bounded on the support, so
$\|D_v(h,0)\psi_h\|=O(1)$. The undeformed square identity of
Proposition~\ref{prop:square}, with the fermion term bounded below by
$-Ch|x|$, gives
\[
 \|p\psi_h\|^2+h^2\sum_{i<j}\|[x_i,x_j]\psi_h\|^2
 \leq C\bigl(\|D_v(h,0)\psi_h\|^2+h\|\psi_h\|^2\bigr).
\]
For \eqref{fast:growing-mass}, average the sixteen squared charges. The
energy is $o(m_h^2)$, and the fermion term on the support has norm at most
$C(h+m_h)$. In each shifted bosonic square use
$|a-b|^2\geq\frac12|a|^2-|b|^2$. The subtracted terms are bounded by
$Cm_h^2\|\psi_h\|^2$. Since $m_h^2=O(h)$, the same conclusion follows.
These form estimates extend from the invariant core by
Lemma~\ref{lem:domains}. Finally, on a compact set separated from the
commuting locus, the continuous function
$\sum_{i<j}|[x_i,x_j]|^2$ has a positive minimum. The second estimate
therefore makes the norm on that set tend to zero.
\end{proof}

Fix a commuting configuration with $k\geq2$ distinct joint centers and
multiplicities $n_1,\ldots,n_k$. Write
\[
 \mathfrak k=\mathfrak k_0\oplus\mathfrak n,
 \qquad
 \mathfrak k_0=
 \left(\bigoplus_{\alpha=1}^k\mathfrak u(n_\alpha)\right)
 \cap\mathfrak{su}(N),
\]
orthogonally for the trace metric. The full block unitary group $K_0=\prod_{\alpha=1}^k\U(n_\alpha)$
acts on this splitting; its central torus is $Z_0$. Its scalar diagonal
subgroup acts trivially by conjugation.
Choose a unit spatial vector $e$ whose scalar products with the centers
are all distinct. Complete $e$ to an orthonormal spatial frame and use
$r$ for the eight remaining frame indices. The definitions of $A_i$,
$B_{ij}$, and $L_i$ are transformed as tensors when this frame is used.

\begin{lemma}[Exact slice and charge]\label{lem:slice}
In a neighborhood of the chosen configuration, gauge-equivariant
functions are equivalent to $K_0$-equivariant functions on
\[
 x=b+z,\qquad b\in\mathfrak k_0^9,
 \qquad z\in e^\perp\otimes\mathfrak n.
\]
The block spectra of $b_e$ lie in fixed disjoint windows. Up to a fixed
positive orbit-volume constant, the density on the slice is
\begin{equation}
 j(b)=|\det_{\mathfrak n}K_e(b)|,
 \qquad K_e(b)=\ad(b_e)|_{\mathfrak n}.
 \label{fast:density}
\end{equation}
It is independent of $z$. The exact normalization is as follows. Give
$\U(N)/K_0$ the invariant volume induced by the trace metric on the
off-block color space, let $V_C$ be its total volume, and use normalized
Haar measure when integrating the gauge orbit. With fixed labeled
spectral windows the integrated slice measure is
$V_Cj(b)\,\dd b\,\dd z$. Thus, with
$d_f=8\dim\mathfrak n$, the unitary flat-measure pullback is
\begin{equation}
 \psi(b,t)=V_C^{1/2}j(b)^{1/2}h^{-d_f/4}
                         \Psi(b+t/\sqrt h).
 \label{fast:unitary-normalization}
\end{equation}
Here the original field is restricted to the slice on the right.
A fixed labeled chart has no additional permutation multiplicity.
When all labels are used, division by the number of permutations of
equal-sized blocks gives the quotient measure used in
Definition~\ref{def:profile-space}. The same positive normalization is
used in the inverse pullback for every test state. In these coordinates,
\begin{equation}
 D_v(h,m)=\sqrt h\,E_v(b)+D_v^s(h,m)+R_v(h,m).
 \label{eq:slice-charge}
\end{equation}
Here $D_v^s$ is precisely the charge on $\mathfrak k_0^9$, including
all its $h[b_i,b_j]$ multipliers. The operator $E_v(b)$ is an odd
linear oscillator in $(p_t,t)$ using only fast fermions, with smooth
coefficients independent of $h,m$. On a smaller chart the raw estimates
imply
\begin{equation}
 \|h^{-1/2}p_b\psi_h\|+\|p_t\psi_h\|+\|t\psi_h\|\leq C.
 \label{fast:chart-raw}
\end{equation}
\end{lemma}

\begin{proof}
The spectral projections of $x_e$ onto the chosen windows identify its
block subspaces uniquely up to $K_0$. Smooth local choices of frames
give the slice. The differential of conjugation in an off-block color
$w\in\mathfrak n$ has missing component $[w,b_e]$. All other differential
components occur below this block in the Jacobian matrix, whose other
diagonal blocks are the identity on the slice and the fixed homogeneous
space measure. This proves \eqref{fast:density} exactly. Invertibility
of $K_e$ and uniform bounds on its inverse follow from the spectral
gaps. The same spectral-window parametrization allows arbitrary $z$;
small $z$ is needed only to stay in the original compact neighborhood.

Here are coordinate formulas that also fix the content of the
remainder. Let $(n_A)$ be an orthonormal basis of $\mathfrak n$ and set
\[
 w_A(b)=-K_e(b)^{-1}n_A,
 \qquad [w_A(b),b_e]=n_A.
\]
Write $\theta_f$ and $\theta_s$ for fast and slow color Clifford
variables. Before density conjugation, gauge equivariance gives
\begin{equation}
 \langle[w,b_e],p_e\rangle
 =-\sum_r\langle[w,b_r+z_r],p_r\rangle
   -\frac{\ii}{2}\theta^t\ad(w)\theta.
 \label{fast:gauge-momentum}
\end{equation}
All coefficients in this identity multiply derivatives on their left.
In particular, for $t_e=0$, the exact fast operator is
\begin{align}
 E_v(b)={}&\sum_r\theta_f(A_r)\cdot p_{t_r}
 -\sum_{A,r}\theta_f^A(A_e)
       \langle[w_A,b_r],p_{t_r}\rangle\notag\\
 &+\sum_{i<j}\theta_f(B_{ij})\cdot
                 \bigl([b_i,t_j]+[t_i,b_j]\bigr).
 \label{fast:E-coordinate}
\end{align}
If $\Pi_0$ and $\Pi_\perp$ denote the color projections onto
$\mathfrak k_0$ and $\mathfrak n$, respectively, the remaining terms
before the density correction are
\begin{align}
 R_v={}&\sum_{i<j}\theta(B_{ij})\cdot[t_i,t_j]
 +\frac m{\sqrt h}\sum_i\theta_f(L_i)\cdot t_i\notag\\
 &-\sum_{A,r}\theta_f^A(A_e)
       \langle\Pi_\perp[w_A,t_r],p_{t_r}\rangle\notag\\
 &-\frac1{\sqrt h}\sum_{A,r}\theta_f^A(A_e)
       \langle\Pi_0[w_A,t_r],p_{b_r}\rangle
 -\frac{\ii}{2}\sum_A\theta_f^A(A_e)
                          \theta^t\ad(w_A)\theta.
 \label{fast:R-coordinate}
\end{align}
The formula after density conjugation is obtained by replacing every
$p_{b_i}$ by $p_{b_i}+\frac{\ii}{2}\partial_{b_i}\log j$, and putting
the added terms in $R_v$. In this spatial frame $j$ depends only on
$b_e$, so the added term is simply
\begin{equation}
 \frac{\ii}{2}\theta_s(A_e)\cdot\nabla_{b_e}\log j.
 \label{fast:density-term}
\end{equation}
These formulas include the Clifford signs of the actual full module;
when it is written as a graded tensor product, the fast odd operator
has the usual slow-parity factor. They show in particular that the
only derivative terms in $R_v$ have types $t p_t$ and
$h^{-1/2}t p_b$. All coefficients and their derivatives are bounded on
compact subcharts.

Tangential derivatives on the linear slice are restrictions of
ambient derivatives. Gauge integration and the bounded positive
density transfer the first raw estimate to these derivatives; density
conjugation adds a bounded multiplier. The substitution $z=t/\sqrt h$
then bounds $p_t$ and $p_b/\sqrt h$. Finally,
\[
 [x_e,x_r]_{\mathfrak n}=[b_e,z_r]
\]
and the uniform inverse of $K_e$ transfer the second raw estimate to
$\|t\psi_h\|\leq C$.
\end{proof}

\subsection{The neutral fast oscillator}

The slow variables and their fermions are fixed by $Z_0$. Thus a
$K_0$-equivariant slice function belongs, in its fast variables and
fast fermions together, to the fixed Hilbert space of $Z_0$-neutral
states. We always restrict $E_v(b)$ to this space. Neutrality means
invariance for the whole central torus; separate invariance in every
individual real color plane is not required.

\begin{lemma}[The neutral fast line]\label{lem:fast-line}
Let $c$ range in a compact set of scalar block configurations with
distinct centers and with the chosen $e$-separations bounded below.
On the neutral space, $E_v(c)$ has a one-dimensional even kernel,
independent of $v$, and a gap bounded below uniformly on this compact
set and on unit spinors $v$. For $b$ sufficiently close to that set,
including noncommuting $b$, it has an even rank-one kernel projection
$P_v(b)$ and a uniformly bounded inverse on $\Ran(1-P_v(b))$.

Let $\mathcal W^s$ be the oscillator Sobolev spaces defined below.
For every prescribed finite integer $L$, the neighborhood can be
chosen so that the parameter derivatives of $P_v$ and a normalized
section $\chi_v(b)$ through order $L$ have bounded oscillator moments
through order $L$. The reduced inverse, extended by zero on the
kernel, and its parameter derivatives through order $L$ are uniformly
bounded as maps $\mathcal W^s\to\mathcal W^{s+1}$ for
$0\leq s\leq L$. The section can be chosen
$K_0$-equivariantly so that its Berry form vanishes at $b=c$ in every
slow direction. At scalar centers these sections are compatible with
block permutations and the $\Spin(3)\times\Spin(6)$ action, with no
additional scalar character.
\end{lemma}

\begin{proof}
We first compute the scalar-center oscillator and its gap. Perturbing
that operator will give the exact fast line for nearby noncommuting
blocks and its regularity. We then choose the phase needed for diagonal
compression and for transport between cluster charts.

\smallskip\noindent\emph{The oscillator at scalar centers.}
Change from the $e$-slice to the perpendicular slice for the
linear conjugation directions $([w,c_i])_i$, including the positive
linear density factor. For one real two-dimensional color plane
joining two distinct blocks, let $\lambda\in\R^9$ be their signed
center difference, $r=|\lambda|$, and $J$ the orthogonal complex
structure on that plane, so that $\ad(c_j)=\lambda_jJ$. There are
sixteen real transverse bosonic coordinates: eight spatial directions
times two color directions. The charge on them is
\[
 \theta_f(A_i)\cdot p_i^q+
 \sum_{i<j}\theta_f(B_{ij})\cdot
                       J(\lambda_iq_j-\lambda_jq_i),
 \qquad q\perp\lambda,\quad p^q\perp\lambda.
\]
To check the square, rotate $\lambda$ to $r e_9$ and put
$s=v^t\gamma^9v$. The summand becomes
\[
 \sum_{i=1}^8\bigl(
 \theta_f(\gamma^iv)\cdot p_i^q
 +r\theta_f(\gamma^i\gamma^9v)\cdot Jq_i\bigr),
\]
and the Clifford identity used in Proposition~\ref{prop:square} gives
\begin{equation}
 \frac12\bigl((p^q)^2+r^2q^2\bigr)
 +rs\,p^q\cdot Jq
 +\frac{\ii r}{2}\theta_f^t
                   \bigl[J\otimes(s\Id-\gamma^9)\bigr]\theta_f.
 \label{fast:one-plane-square}
\end{equation}
There are no cross terms between disjoint Clifford summands. Summing
their orbital terms before imposing neutrality gives the gauge
generator for the single central color
$\sum_j(v^t\gamma^jv)c_j$. Equation~\eqref{fast:gauge-momentum} for
this central generator cancels the scalar-spin terms in
\eqref{fast:one-plane-square}. Hence, on the neutral space,
\begin{equation}
 E_v(c)^2=H_{\mathrm{osc}}(c)
       -\frac{\ii}{2}\theta_f^t
                   \bigl(\ad(c_j)|_{\mathfrak n}\otimes\gamma^j\bigr)
                    \theta_f.
 \label{fast:neutral-square}
\end{equation}
In particular this square is independent of $v$.

Each real color plane contributes bosonic ground energy $8r$.
Its thirty-two real fermions give sixteen Clifford oscillators of
frequency $r$, with unique spin ground energy $-8r$. In the
color-complex polarization the latter fills eight modes and leaves
eight empty, so its central color charge is zero; the positive
Gaussian is neutral as well. Their product is therefore in the
neutral space. The nonnegative oscillator occupation numbers show
that this product is the only zero vector up to scalar, and every
other eigenvalue of the square is bounded below by the smallest
center distance. Restriction to the neutral space cannot reduce
that bound. This proves the kernel and gap assertions at centers.

\smallskip\noindent\emph{Perturbation to noncommuting blocks.}
The subsequent first-order estimates use both differentiated projections
and oscillator moments. At a reference center $c_*$ define
$\mathcal W^s$ by the norm
$\|(1+\mathcal N)^{s/2}f\|$, where $\mathcal N$ is its positive
bosonic oscillator number operator; the fermion space is finite.
These are restricted to neutral vectors. The creation and
annihilation relations and \eqref{fast:neutral-square} imply
\[
 \|p_tf\|+\|tf\|\leq C(\|E_v(c_*)f\|+\|f\|).
\]
They also imply, on a circle enclosing zero alone, that
$(E_v(c_*)-\zeta)^{-1}:\mathcal W^s\to\mathcal W^{s+1}$ is bounded
for every fixed $s$. For completeness, the square is an oscillator
number operator plus a bounded spin matrix. Each application of
$t$ or $p_t$ changes the number by one, and its coefficient is bounded
by a constant times $(1+\mathcal N)^{1/2}$. The resolvent can be
written $(E_v(c_*)+\zeta)(E_v(c_*)^2-\zeta^2)^{-1}$; its two factors
respectively lose one and gain two such orders. This proves the
claimed bound, including nonintegral $s$ if needed.

The central torus fixes $b$, and $K_e(b)$, its inverse, and all
$\ad(b_r)|_{\mathfrak n}$ commute with its action. Thus $E_v(b)$ preserves
the same neutral space even when $b$ is noncommuting. Its real linear
coefficients make it symmetric on the neutral Schwartz core.
The coefficient formula \eqref{fast:E-coordinate} gives, for every
fixed finite range of $s$,
\[
 \|E_v(b)-E_v(c_*)\|_{\mathcal W^{s+1}\to\mathcal W^s}
 \leq C_s|b-c_*|.
\]
Choose the neighborhood so that the resolvent Neumann series
converges on each of these finitely many spaces. On $\mathcal W^0$
the same relative bound, made smaller than one, gives self-adjointness
on the reference domain. The contour projection has constant rank
one. It commutes with parity, and its range remains entirely even by
continuity. Since $E_v(b)$ preserves this range and reverses parity,
its restriction to the range is zero. The complementary
inverse can equivalently be formed from
$(E_v(b)+P_v(b))^{-1}(1-P_v(b))$. Differentiating its resolvent
formula and using
\[
 \partial_b(E_v(b)-\zeta)^{-1}
 =-(E_v(b)-\zeta)^{-1}(\partial_bE_v(b))
                         (E_v(b)-\zeta)^{-1}
\]
proves all stated finite-order bounds after increasing the finite
range of oscillator orders. None of this uses $[b_i,b_j]=0$.

\smallskip\noindent\emph{A compatible phase.}
For a fixed oriented color pair, the
spin ground line over $\lambda/r\in S^8$ is homogeneous under
$\Spin(9)$. The underlying orthogonal action lifts to the Clifford module
as an honest action because $\Spin(9)$ is simply connected.
Its stabilizer $\Spin(8)$ acts on the line by a character.
A connected semisimple compact group has no nontrivial continuous
characters, so it fixes a chosen unit vector. Transport therefore
defines a $\Spin(9)$-equivariant unit section. Its Berry one-form is
invariant on $S^8$, and hence zero: the tangent representation of
$\Spin(8)$ has no fixed covector. The normalized real positive
Gaussian also has zero Berry form.

For directions in the last six axes, the unit skew ground-frequency
matrix anticommutes with the complex structure $T$ used for the
exterior module. The projection between the corresponding complex
annihilator spaces has squared norm $1/2$ on unit vectors, and is
invertible. The new annihilators are consequently a graph over the
old ones. Their vacuum has a nonzero degree-zero coefficient: if the
lowest nonzero exterior component had positive degree, the graph
annihilation equations in one lower degree would make all ordinary
contractions of that component zero, an impossibility. This graph
vacuum is even. We fix its phase by positive overlap with the empty
$T$-vacuum. The overlap convention is consistent as the last-six
direction varies, by $\Spin(6)$ covariance.

The block group and its block permutations preserve the empty
$T$-vacuum by their natural exterior actions. Thus on last-six center
configurations the product ground section is block invariant and
has exactly the permutation transformation law of the underlying
coordinates, without a scalar factor. Reordering its pair factors
adds no sign, since each is even. These sections are parallel as
functions of the labeled distinct centers. The quotient of two
sections related by a permutation is therefore locally constant.
The space of distinct labeled points in $\R^9$ with trace-zero
weighted mean is connected; removing the coincidence subspaces of
codimension nine does not disconnect it. Agreement on last-six
configurations proves agreement everywhere. The same construction
gives covariance under $\Spin(3)\times\Spin(6)$. These statements
concern the fast line; the graded permutation on slow factors is
still the ordinary one.

Let $c=c(b)$ be the scalar block-center projection. If $\xi(c)$ is
the unit section just constructed, set
\begin{equation}
 \chi_v(b)=\frac{P_v(b)\xi(c(b))}
                     {\|P_v(b)\xi(c(b))\|}.
 \label{fast:section-choice}
\end{equation}
The denominator is bounded away from zero on a smaller neighborhood.
This prescription is $K_0$-equivariant. In internal directions at
$b=c$, $\xi$ is constant and $P(\partial P)P=0$ shows that
$\langle\chi_v,\partial_b\chi_v\rangle=0$. In center directions the
same assertion is the zero Berry form already proved. Smooth
transverse linear coordinate changes, implemented with the positive
square root of the absolute Jacobian, do not change it: they send
the normalized positive real Gaussian to a normalized positive real
Gaussian, whose inner product with its derivative is zero.
\end{proof}

\subsection{Projection and the effective slow charge}

First-order oscillator projection and the cancellation between gauge
spin and slice density occur in the regular Cartan analysis of
Hasler--Hoppe~\cite{HaslerHoppeFactorization2002}. We need the following
uniform estimates also for nearby noncommuting blocks and for the
separate growing-mass regime.

\begin{lemma}[Projection and the slow graph bound]
\label{lem:fast-projection}
In a compact subchart of Lemma~\ref{lem:fast-line}, in either of the
two local regimes,
\begin{equation}
 \|(1-P_v)\psi_h\|\longrightarrow0,
 \qquad
 \|D_v^s(h,m)P_v\psi_h\|
 \leq
 \begin{cases}
 O(1),&m\text{ fixed},\\
 o(m),&\text{in \eqref{fast:growing-mass}}.
 \end{cases}
 \label{eq:fast-projection}
\end{equation}
The second norm is on the full $(b,t)$ space. In the fixed-mass case,
writing $P_v\psi_h=\chi_v(b)\eta_h(b)$ gives a bounded graph norm
for $D_v^s(h,0)\eta_h$, where $\eta_h$ has the ordinary slow
$K_0$-equivariance.
\end{lemma}

\begin{proof}
For each fixed $h$, local ellipticity of the original charge puts a
compactly supported graph-domain vector in $H^1_{\mathrm{loc}}$.
Gauge equivariance and the smooth orbit--slice coordinates transfer
this regularity to the slice, by orbit integration, on compact subcharts. After a compact fast cutoff, $E_v\psi$ and $D_v^s\psi$
are therefore separately in $L^2$: the slow multipliers are bounded
on the compact slow support for that fixed $h$. Local smoothing
justifies the following squared-norm identities before any uniform
bound on these separate operators is invoked.

Choose an even scalar cutoff $\rho_M(t)$ equal to one for $|t|\leq M$,
supported in $|t|\leq2M$, and with gradient bounded by $C/M$. It is
$K_0$-invariant. Equation~\eqref{fast:chart-raw} gives
$\|(1-\rho_M)\psi_h\|\leq C/M$. For fixed $M$, the explicit remainder
formula shows
\[
 \|h^{-1/2}R_v\rho_M\psi_h\|\longrightarrow0.
\]
Indeed the respective bounds are constants depending on $M$ times
$h^{-1/2}$ for $t^2$, spin, and density terms;
$h^{-1/2}\|p_t\psi_h\|$ for $tp_t$;
$h^{-1}\|p_b\psi_h\|$ for $tp_b/\sqrt h$; and $m/h$ for the mass
term. The cutoff commutator with $D_v/\sqrt h$ is bounded by
$C/M+o(1)$. Also $\|D_v\psi_h\|/\sqrt h\to0$ in both regimes.
Consequently
\begin{equation}
 \limsup_{h\to\infty}
 \|(E_v+h^{-1/2}D_v^s)\rho_M\psi_h\|\leq C/M.
 \label{fast:cut-charge}
\end{equation}

Every slow bracket and mass Clifford multiplier anticommutes with
$E_v$. In particular, the potentially large multiplier $h[b,b]$
produces no anticommutator. Only slow differentiation of its
coefficients remains:
\[
 \{E_v,D_v^s\}
 =-\ii\sum_{i,Q}\theta_s^Q(A_i)\,\partial_{b_i^Q}E_v.
\]
On $\rho_M\psi_h$ its expectation is bounded uniformly for fixed
$M$, because $\partial_bE_v$ is linear in $(p_t,t)$ and
\eqref{fast:chart-raw} bounds their norms. Squaring
\eqref{fast:cut-charge}, dropping the nonnegative slow square, and
using the fast gap thus gives
\[
 \limsup_{h\to\infty}\|(1-P_v)\rho_M\psi_h\|\leq C'/M.
\]
The tail estimate and then $M\to\infty$ prove the first assertion.

For the second assertion, it is enough to estimate the projected
remainder. The rank-one projection, its oscillator moments, and
integration by parts in $t$ give the uniform bound
\begin{equation}
 \|P_vR_v\psi_h\|
 \leq C\left((1+m/\sqrt h)\|\psi_h\|
                       +\|p_b\psi_h/\sqrt h\|\right).
 \label{fast:projected-remainder}
\end{equation}
To see each term, $t^2$ moves onto the oscillator bra; in $tp_t$ the
derivative moves onto that bra and $t$; spin and density terms are
bounded; $tp_b/\sqrt h$ uses the bounded bra $t\chi_v$ and the last
norm on the right; and the fast mass term uses $t\chi_v$ and
$m/\sqrt h$. Thus no norm estimate for the unprojected $R_v\psi_h$
is asserted or needed. Moreover
\[
 [D_v^s,P_v]=-\ii\sum_{i,Q}\theta_s^Q(A_i)
                                      \partial_{b_i^Q}P_v
\]
is bounded: the even projection commutes with all slow Clifford
multipliers. Projecting \eqref{eq:slice-charge} now gives
\[
 D_v^sP_v\psi_h=P_vD_v\psi_h-P_vR_v\psi_h
                                      +[D_v^s,P_v]\psi_h.
\]
Equations~\eqref{fast:chart-raw} and
\eqref{fast:projected-remainder} prove the claimed estimates. Finally,
differentiating $\chi_v\eta_h$ costs a bounded derivative of $\chi_v$;
the fixed mass multiplier is bounded on the chart. Removing the
section therefore gives the asserted undeformed slow graph bound.
Its equivariance follows from \eqref{fast:section-choice}.
\end{proof}

\begin{lemma}[Diagonal compression at scalar blocks]
\label{lem:diagonal-compression}
For fixed $m$, identify the range of $P_v(b)$ with the slow Clifford
fiber by $a\mapsto\chi_v(b)a$. At $b=c$, the diagonal compression of
\eqref{eq:slice-charge} has exactly the slow derivative symbol and
the slow order-zero terms of $D_v^s(h,m)$, apart from terms whose
coefficients tend to zero as $h\to\infty$. More precisely, after
discarding explicitly $h^{-1/2}$-small coefficients, all additional
order-zero coefficients are smooth and vanish at $b=c$.
\end{lemma}

\begin{proof}
The Berry term vanishes by Lemma~\ref{lem:fast-line}. The outer fast
fermion in each $tp_t$ or $tp_b$ term makes its diagonal matrix
element zero, also when a slow derivative acts on the even section.
The slow-color part of $[t_i,t_j]$ has zero Gaussian expectation:
on each cross-block color space the covariance is a symmetric
spatial matrix times the color identity, whose contraction with the
antisymmetric color bracket is zero. Its fast-color part and the
other purely fast odd terms have zero matrix element by parity.

We verify the remaining orbit-spin and density cancellation for
every slow color, including internal traceless colors. For a slow
orthonormal basis color $q\in\mathfrak k_0$, put
\[
 K_q=\ad(q)|_{\mathfrak n},\qquad K_e=\ad(c_e)|_{\mathfrak n}.
\]
There can be no term containing three slow fermions in the
orbit-spin substitution: for $w\in\mathfrak n$ the slow-to-slow
block of $\ad(w)$ is zero. Its fast-to-fast block, after the outer
fast fermion is included, has odd fast degree and has zero
expectation. Only the mixed slow-fast block contributes.

On a cross-block space joining blocks $\alpha,\beta$, one has
$K_e=\delta_{\alpha\beta,e}J_{\alpha\beta}$ with
$\delta_{\alpha\beta,e}\ne0$. The operator $K_q$ preserves this
space and is real skew. The fast ground two-point matrix is
$\frac12\Id$ plus an imaginary skew matrix whose color factor is
$J_{\alpha\beta}$ and whose spin factor is the signed unit
center-difference Clifford matrix. Its imaginary part in the mixed
contraction has color trace
\begin{equation}
 \Tr\bigl(J_{\alpha\beta}K_qK_e^{-1}\bigr)
 =\delta_{\alpha\beta,e}^{-1}\Tr K_q=0.
 \label{fast:imaginary-cancellation}
\end{equation}
The equality follows by cyclicity and $J_{\alpha\beta}^2=-\Id$;
it does not require $q$ to be central. The spin coefficient can
therefore be left arbitrary in this cancellation.

For clarity, the real scalar contraction can be computed without a
symmetry argument. With the fast fermion ordered before the slow
one, the mixed entries of $\ad(w_A)$ are
$[w_A,q]=-K_qw_A$, and their two orders combine to twice this
coefficient. Multiplying by
$-\frac{\ii}{2}\theta_f^A(A_e)$ from
\eqref{fast:R-coordinate} and using
$\langle\theta_{f,\mu}^A\theta_{f,\nu}^B\rangle_{\mathrm{scalar}}
=\frac12\delta_{AB}\delta_{\mu\nu}$ gives
\begin{equation}
 -\frac{\ii}{2}\Tr_{\mathfrak n}(K_e^{-1}K_q)
                                      \theta_s^q(A_e).
 \label{fast:scalar-spin}
\end{equation}
On the other hand, differentiating the determinant in
\eqref{fast:density} gives
\[
 \partial_{b_e^q}\log j\big|_{b=c}
                  =\Tr_{\mathfrak n}(K_e^{-1}K_q).
\]
The density term \eqref{fast:density-term} is exactly the opposite
of \eqref{fast:scalar-spin}. This proves the cancellation for every
$q$. There are no terms to check in other spatial directions:
$j$ depends only on $b_e$, and only the $e$-momentum was eliminated.
The finitely many remaining $h^{-1/2}$ coefficients are explicitly
visible in \eqref{fast:R-coordinate}. Smoothness and the oscillator
bounds prove the last assertion of the lemma.
\end{proof}

Fast projection now locates the leading state and identifies its diagonal
slow charge. The next section cancels the remaining transverse charge
output, so that the slow equation can be passed to limits by duality.
